\part{確率論}
	\chapter{表記}
		\begin{itemize}
			\item 特に断らない限り、問題は確率空間$(\Omega,\mathcal{F},P)$で議論されているものと約束する。
		\end{itemize}
	\chapter{独立性に関する諸定理}
	\section{\texorpdfstring{$A_1,\dotsc,A_n$}{A\_1,...,A\_n}が独立ならば、そのうちどれか1つ\texorpdfstring{$A_l$}{A\_l}を\texorpdfstring{$A_l^\complement$}{A\_l\string^∁}に置き換えたものも独立である}
		\begin{proof}
			\begin{align*}
				\func{P}{A_l^\complement\cap\left(\bigcap_{i\neq l}^n A_i\right)} &= \func{P}{\bigcap_{i\neq l}^n A_i} - \func{P}{A_l\cap\left(\bigcap_{i\neq l}^n A_i\right)} = \func{P}{\bigcap_{i\neq l}^n A_i} - \func{P}{\bigcap_{i=1}^n A_i} \\
				&= \prod_{i\neq l}^n P(A_i) - \prod_{i=1}^n P(A_i) \quad (\because\;\text{独立性の仮定}) \\
				&= (1-P(A_l))\prod_{i\neq l}^n P(A_i) = P(A_l^\complement)\prod_{i\neq l}^n P(A_i)
			\end{align*}
		\end{proof}
	\section{\texorpdfstring{$A_1,\dotsc,A_n$}{A\_1,...,A\_n}が独立ならば\texorpdfstring{$A_1^\complement,\dotsc,A_n^\complement$}{A\_1\string^∁, ..., A\_n\string^∁}も独立である}
		\begin{proof}
			直前の定理を繰り返し用いる。
		\end{proof}
		\begin{proof}
			(別証)
			\par
			包除原理を用いる。
			\begin{align*}
				\func{P}{\bigcap_{i=1}^n A_i^\complement} &= 1 - \func{P}{\left(\bigcap_{i=1}^n A_i^\complement\right)^\complement} = 1 - \func{P}{\bigcup_{i=1}^n A_i} = 1 - \sum_{j=1}^n (-1)^{j-1}\sum_{I\in\combi{\{1,\dotsc,n\}}{j}} \func{P}{\bigcap_{i\in I}A_i} \\
				&= 1 + \sum_{j=1}^n (-1)^j\sum_{I\in\combi{\{1,\dotsc,n\}}{j}} \prod_{i\in I}P(A_i) = 1 + \sum_{I\in \text{powset}(\{1,\dotsc,n\})\setminus\emptyset} (-1)^{|I|}\prod_{i\in I}P(A_i) \\
				&= \prod_{i=1}^n (1-P(A_i)) = \prod_{i=1}^n P(A_i^\complement)
			\end{align*}
		\end{proof}
	\section{\texorpdfstring{$A_1,\dotsc,A_n$}{A\_1,...,A\_n}が独立ならば、そのうちどれか2つ\texorpdfstring{$A_p,A_q$}{A\_p, A\_q}を合併したものも独立である。}
		\begin{proof}
			\quad\par
			\begin{align*}
				&\quad \func{P}{(A_p\cup A_q)\cup\left(\bigcap_{i\neq p,q}^n A_i\right)} \\
				&= \func{P}{\left(\left(\bigcap_{i\neq p,q}^n A_i\right)\cap A_p\right)\cup\left(\left(\bigcap_{i\neq p,q}^n A_i\right)\cap A_q\right)} \\
				&= \func{P}{\left(\left(\bigcap_{i\neq p,q}^n A_i\right)\cap A_p\right)} + \func{P}{\left(\left(\bigcap_{i\neq p,q}^n A_i\right)\cap A_q\right)} \\
				&\quad - \func{P}{\left(\left(\bigcap_{i\neq p,q}^n A_i\right)\cap A_p\right)\cap\left(\left(\bigcap_{i\neq p,q}^n A_i\right)\cap A_q\right)} \\
				&= \func{P}{\bigcap_{i\neq p,q}^n A_i}P(A_p) + \func{P}{\bigcap_{i\neq p,q}^n A_i}P(A_q) \\
				&\quad - \func{P}{\left(\bigcap_{i\neq p,q}^n A_i\right)\cap(A_p\cap A_q)} \\
				&= \func{P}{\bigcap_{i\neq p,q}^n A_i}P(A_p) + \func{P}{\bigcap_{i\neq p,q}^n A_i}P(A_q) \\
				&\quad - \func{P}{\bigcap_{i\neq p,q}^n A_i}P(A_p)P(A_q) \\
				&= (P(A_p) + P(A_q) - P(A_p)P(A_q))\func{P}{\bigcap_{i\neq p,q}^n A_i} \\
				&= P(A_p\cup A_q)\func{P}{\bigcap_{i\neq p,q}^n A_i}
			\end{align*}
		\end{proof}
	\section{\texorpdfstring{$A_1,\dotsc,A_n$}{A\_1,...,A\_n}が独立ならば、\texorpdfstring{$I_1\sqcup I_2 = \{1:n\}$}{I\_1 ⊔ I\_2}なる\texorpdfstring{$I_1,I_2$}{I\_1, I\_2}に対して\texorpdfstring{$\bigcup_{i\in I_1}A_i$}{∪\_\{i∈I\_1\}A\_i}と\texorpdfstring{$\bigcup_{j\in I_2}A_j$}{∪\_\{j∈I\_2\}A\_j}とは独立である。}
		\begin{proof}
			直前の定理を繰り返し用いる。
		\end{proof}